\documentclass[a4paper,11pt]{article}

\usepackage{revy}
\usepackage[utf8]{inputenc}
\usepackage[T1]{fontenc}
\usepackage[danish]{babel}


\revyname{MatematikRevyen}
\revyyear{2026}
\version{1}
\eta{$2$ minutter}
\status{Færdig}

\title{Tællelig}
\author{Julius Pallesen '24}
\melody{Madame Monsieur: ``Mercy'' (forkortet)}

\begin{document}
\maketitle

\begin{roles}
\role{S}[Sanger] Sanger
\end{roles}

\begin{props}
    \prop{Pose med kridt} Pose med uendelig mange kridt. Det skal forestille at være tavlekridt, men i virkeligheden er det skolekridt så de kan deles ud til publikum
\end{props}

\begin{song}
    \scene{S har en pose med uendelig mange skolekridt og heller en god del ud på gulvet. Bandet begynder, og S begynder at tælle kridt}

\sings{S} 1, 2, 3, – 4, 5, 6
7, 8, 9, 10
Jeg kunne blive ved
til uendelig, plus 10

Jeg så træt af at tæl' til uendelig
Det tager fucking lang tid, at gøre det
Åh ja, at tæl' til uendelig
Der så fucking mange tal, man skal tælle til
Og det er slet ikke det, som jeg gerne vil
Åh men, nogen skal jo gøre det

Men mængden er diskret, den er tæl - lelig
Det faktisk ik' så mang', når man sam – ligner det
De fleste mængder er faktisk o - vertæl'lig'
Men den her pose kridt, den er – tæl'lig

Tællelig, tællelig, tællelig, tællelig \act{Begynder at dele ud til publikum}
Tællelig, tællelig, ik' overtællelig
Tællelig, tællelig, tællelig, tællelig
Tællelig, tællelig, ik' overtællelig
Tællelig, tællelig, tællelig, tællelig
Tællelig, tællelig, ik' overtællelig

\end{song}

\end{document}

